\scanpage{541}%
\mspage{452}%
% archive sheet number: 540
est isomorphe \unsure{:|à} un syst. inductif constant, de valeur $z$. En choisissant un tel isom

\[z\mathrel{\overset{\overline\alpha_j}{\overset{\sim}{\relbar}}}x_{\unsure{j|i}}^\natural,\]

on trouve $z\xrightarrow{\alpha}\underline{x}$, et le système des $p_i$ définit des

\diagram{\begin{tikzcd}[ampersand replacement=\&] x_{\unsure{j|i}}\arrow[r,"{\overline p_i}"]\arrow[rr,bend right=22,dash,"{\beta_i\overset{\mathrm{déf}}{=}\overline\alpha_j^{-1}\,\overline p_i}"'] \& x_i^\natural\arrow[r,"{\overline\alpha_j^{-1}}"] \& \deleted{\unsure{x|z}}z \end{tikzcd}},

et on voit \deleted{\unsure{aussitôt}} \unsure{facilemt|facilement} que l'on aura

\[\beta\alpha=\mathrm{id}_z,\qquad\alpha\beta=\mathrm{id}_{\underline{x}},\qquad\text{qed.}\]

Je \uline{\unsure{peux} maintenant prouver 4.15.3.} C'est en effet la conjonction de 4.15.7.5 (p. 448), et de 4.8.13, qui nous dit \inserted{\unsure{en.|env.}} que si $C$ est \uline{\unsure{Karoubienne|Karoubien}} \unsure{$\nearrow$|*} (condition vérifiée si $C$ est sans projecteurs), alors tt syst. inductif dans $C$, \unsure{grd|petit|pt} devant $\pi_0$, où $\pi_0\geq\operatorname{card}\operatorname{\unsure{red|réd|nerf}}C$, est essentiellement constant.

Qed.
